\subsection{Cylinder of a map}

Let X and Y be topological spaces, and let f be a continuous map from X to Y. Denote by U the topological sum space of the space \(U_{1} = X \times I\) and the space \(U_{2} = Y\), and identify \(U_{1}\) and \(U_{2}\) with subspaces of U via the canonical injections. Let R be the finest equivalence relation on U for which the points \((x,1)\) of \(U_{1}\) and \(f(x)\) of \(U_{2}\) are equivalent, for every \(x \in X\).

DEFINITION 7. --- We call the cylinder of the map f the quotient topological space U/R and denote it by \(\mathrm{Cyl}(f)\).

Denote by \(\alpha_{f}\colon\mathrm{X}\times\mathbf{I}\to\mathrm{Cyl}(f)\) and \(\beta_{f}\colon\mathrm{Y}\to\mathrm{Cyl}(f)\) the restrictions to \(\mathrm{U}_{1}\) and \(\mathrm{U}_{2}\) of the canonical surjection from U onto U/R. The map \(\alpha_{f}\) is a homotopy and its term is \(\beta_{f}\circ f\).

PROPOSITION 4 (Universal property of cylinders)

Let Z be a topological space, let \(\beta: \mathrm{Y} \to \mathrm{Z}\) be a continuous map and let \(\alpha: \mathrm{X} \times \mathbf{I} \to \mathrm{Z}\) be a homotopy whose term is \(\beta \circ f\). There exists a unique continuous map \(\varphi\) from \(\operatorname{Cyl}(f)\) to Z such that \(\alpha = \varphi \circ \alpha_f\) and \(\beta = \varphi \circ \beta_f\).

The map \(\psi\) from U to Z which coincides with \(\alpha\) on \(\times\) and with \(\beta\) on Y is continuous. We have \(\alpha(x,1)=\beta(f(x))\) for all \(x\in X\), therefore \(\psi\) passes to the quotient to define a continuous map \(\varphi\) from \(\mathrm{Cyl}(f)\) to Z such that \(\alpha=\varphi\circ\alpha_{f}\) and \(\beta=\varphi\circ\beta_{f}\). Since the images of \(\alpha_{f}\) and \(\beta_{f}\) cover \(\mathrm{Cyl}(f)\), \(\varphi\) is the only map from \(\mathrm{Cyl}(f)\) to Z satisfying these relations.

Example 1. --- Let X' and Y' be topological spaces and \(f': \mathrm{X}' \to \mathrm{Y}'\) a continuous map. Let \(u: \mathrm{X} \to \mathrm{X}'\) and \(v: \mathrm{Y} \to \mathrm{Y}'\) be continuous maps such that \(f' \circ u = v \circ f\). There exists a unique continuous map \(w: \operatorname{Cyl}(f) \to \operatorname{Cyl}(f')\) such that \(w \circ \alpha_f = \alpha_{f'} \circ (u \times \operatorname{Id}_\mathbf{I})\) and \(w \circ \beta_f = \beta_{f'} \circ v\).

By prop. 4, applied to Z=Y and \(\beta = \operatorname{Id}_{Y}\), there exists a unique continuous map \(\gamma_{f}\colon\operatorname{Cyl}(f)\to\operatorname{Y}\) such that \(\gamma_{f}(\alpha_{f}(x,s))=f(x)\) and \(\gamma_{f}(\beta_{f}(y))=y\) for \(x\in X\), \(s\in I\), and \(y\in Y\).

PROPOSITION 5. --- The map \(\alpha_{f}\) induces a homeomorphism of X × {[}0,1{[} onto an open subset of \(\mathrm{Cyl}(f)\). The map \(\beta_{f}\) defines a homeomorphism of Y onto the complement of this open set. The map \(\beta_{f} \circ \gamma_{f}\) is a continuous retraction of \(\mathrm{Cyl}(f)\) onto \(\beta_{f}(\mathrm{Y})\).

The set X × {[}0,1{[} is an open subset of U, saturated for the relation R, and the equivalence relation induced by R in X × {[}0,1{[} is the equality relation. The first assertion follows from this (TG, I, p. 23, cor. 10).

The complement of \(\alpha_{f}(\mathrm{X}\times[0,1])\) is \(\beta_{f}(\mathrm{Y})\). Since one has \(\gamma_{f}\circ\beta_{f}=\mathrm{Id}_{\mathrm{Y}}\), \(\beta_{f}\) defines a homeomorphism of Y onto \(\beta_{f}(\mathrm{Y})\) and \(\beta_{f}\circ\gamma_{f}\) is a continuous retraction of the canonical injection of \(\beta_{f}(\mathrm{Y})\) into \(\mathrm{Cyl}(f)\).

The closed subspace \(\beta_{f}(\mathrm{Y})\) of \(\operatorname{Cyl}(f)\) is called the base of the cylinder of \(f\). The map \(\beta_{f} \circ \gamma_{f}\) is called the canonical retraction of \(\operatorname{Cyl}(f)\) onto its base.

Consider the maps \(\sigma_{1}\colon\mathrm{U}_{1}\times\mathbf{I}\to\mathrm{Cyl}(f)\) and \(\sigma_{2}\colon\mathrm{U}_{2}\times\mathbf{I}\to\mathrm{Cyl}(f)\) defined by

\[
\begin{array}{c c} \sigma_ {1} ((x, s), t) = \alpha_ {f} (x, (1 - t) s + t) & \text {for} (x, s) \in \mathrm{X} \times \mathbf {I} \text {and} t \in \mathbf {I} \\ \sigma_ {2} (y, t) = \beta_ {f} (y) & \text {for} y \in \mathrm{Y} \text {and} t \in \mathbf {I}. \end{array}
\]

They are continuous. For \(x \in X\) and \(t \in I\) , we have

\[
\sigma_ {1} ((x, 1), t) = \alpha_ {f} (x, 1) = \beta_ {f} (f (x)) = \sigma_ {2} (f (x), t).
\]

There therefore exists a unique map \(\sigma_{f}\) from \(\mathrm{Cyl}(f)\times\mathbf{I}\) to \(\mathrm{Cyl}(f)\) such that \(\sigma_{f}\circ(\alpha_{f}\times\mathrm{Id}_{\mathbf{I}})=\sigma_{1}\) and \(\sigma_{f}\circ(\beta_{f}\times\mathrm{Id}_{\mathbf{I}})=\sigma_{2}\). The map \(\sigma_{f}\) is continuous (I, p. 19, prop. 10).

Proposition 6. — The map \(\sigma_{f}\colon\mathrm{Cyl}(f)\times\mathbf{I}\to\mathrm{Cyl}(f)\) is a strong contraction of \(\mathrm{Cyl}(f)\) onto its base. Its term is the canonical retraction of \(\mathrm{Cyl}(f)\) onto its base.

The map \(\sigma_f\) is a homotopy. The relations \(\sigma_f(\alpha_f(x,s),0) = \alpha_f(x,s)\) and \(\sigma_f(\beta_f(y),0) = \beta_f(y)\) imply that the origin of \(\sigma_f\) is the identity map of \(\mathrm{Cyl}(f)\). Denote by \(r_f\) the term of \(\sigma_f\). The relations

\[
\sigma_ {f} \left(\alpha_ {f} (x, s), 1\right) = \alpha_ {f} (x, 1) = \beta_ {f} (f (x)) = \left(\beta_ {f} \circ \gamma_ {f}\right) \left(\alpha_ {f} (x, s)\right)
\]

and \(\sigma_{f}(\beta_{f}(y),1) = \beta_{f}(y) = (\beta_{f}\circ \gamma_{f})(\beta_{f}(y))\) imply that \(r_f\) is the canonical retraction \(\beta_{f}\circ \gamma_{f}\) of \(\mathrm{Cyl}(f)\) onto its base.

For \((x,s)\in \mathbf{X}\times \mathbf{I}\) and \(t\in \mathbf{I}\) , we have

\[
r _ {f} \left(\sigma_ {f} \left(\alpha_ {f} (x, s), t\right)\right) = r _ {f} \left(\alpha_ {f} (x, s (1 - t) + t)\right) = \beta_ {f} (f (x)) = r _ {f} \left(\alpha_ {f} (x, s)\right).
\]

For \(y \in Y\) and \(t \in I\), we have \(r_{f}(\sigma_{f}(\beta_{f}(y), t)) = r_{f}(\beta_{f}(y))\). Consequently, \(\sigma_{f}\) is a strong contraction of \(\operatorname{Cyl}(f)\) onto its base.

The map \(\sigma_{f}\) is called the canonical contraction of the cylinder of \(f\) onto its base.

Remarks. --- 1) Let A be a subset of X × I and let A₁ be the set of points x ∈ X such that (x,1) ∈ A. Then we have \(\bar{\alpha}_{f}^{-1}(\alpha_{f}(A)) = A \cup f^{-1}(f(A_{1})) \times \{1\}\) and \(\beta_{f}^{-1}(\alpha_{f}(A)) = f(A_{1})\). Consequently, the map \(\alpha_{f}\) is closed (resp. open) if f is.

\begin{enumerate}
\def\labelenumi{\arabic{enumi})}
\setcounter{enumi}{1}
\tightlist
\item
  Suppose that the map \(f\) is proper. Let \(P\) be a point of \(\operatorname{Cyl}(f)\). If \(P = \alpha_f(x,t)\), with \(0 \leqslant t < 1\) and \(x \in X\), we have \(\overline{\alpha}_f^1(P) = (x,t)\). In the contrary case, there exists \(y \in Y\) such that \(P = \beta_f(y)\)
\end{enumerate}

and \(\bar{\alpha}_{f}^{1}(P) = \bar{f}^{1}(y) \times \{1\}\). This proves that the fibers of the map \(\alpha_{f}\) are quasi-compact. The map \(\alpha_{f}\) is then proper (TG, I, p. 75, th. 1), since it is closed.

By prop. 5 and TG, I, p. 72, prop. 2, the map \(\beta_f\) is itself proper. Consequently, if f is proper, the canonical surjection from U onto \(\mathrm{Cyl}(f)\) is proper, hence, in particular, universally strict (I, p. 20, corollary).

\begin{enumerate}
\def\labelenumi{\arabic{enumi})}
\setcounter{enumi}{2}
\tightlist
\item
  If the spaces X and Y are separated, then so is the cylinder of f. Indeed, let z and \(z'\) be two distinct points of \(\mathrm{Cyl}(f)\) and let us prove that they possess disjoint neighborhoods. We distinguish three cases.
\end{enumerate}

- There exists \((x,t)\) and \((x',t') \in \mathrm{X} \times [0,1[\) such that \(z = \alpha_f(x,t)\) and \(z' = \alpha_f(x',t')\) .

In this case, the assertion follows from the fact that the space \(X \times [0,1]\) is separated (TG, I, p. 54, prop. 7) and that the map \(\alpha_{f}\) induces a homeomorphism from this space onto an open subspace of \(\operatorname{Cyl}(f)\).

- There exists \((x,t)\in \mathrm{X}\times [0,1[\) such that \(z = \alpha_{f}(x,t)\) and \(y^{\prime}\in \mathrm{Y}\) such that \(z^{\prime} = \beta_{f}(y^{\prime})\).

Then, \(\alpha_{f}(\mathrm{X}\times [0,\frac{t + 1}{2}])\) and \(\operatorname {Cyl}(f)\mathbf{-}\alpha_{f}(\mathrm{X}\times [0,\frac{t + 1}{2}])\) are disjoint open neighborhoods of \(z\) and \(z^{\prime}\) in \(\operatorname {Cyl}(f)\).

- There exist \(y\) and \(y'\in Y\) such that \(z=\beta_f(y)\) and \(z'=\beta_f(y')\).

In this case, \(y\neq y'\); since Y is separated, there exists an open neighborhood V of y in Y and an open neighborhood \(V'\) of \(y'\) in Y such that \(V\cap V'=\varnothing\). Then, \((\beta_f\circ\gamma_f)^{-1}(V)\) and \((\beta_f\circ\gamma_f)^{-1}(V')\) are disjoint open subsets of \(\operatorname{Cyl}(f)\) containing respectively \(\beta_f(y)\) and \(\beta_f(y')\).

